% =====================================================================
%  TYPE: Postoperative instructions -- craniotomy / stereotactic biopsy
%  Sections: patient block, instructions, follow-up, emergency
%  box, when to call, medication page. Based on Panagidis 2025-10 and Georgiadis 2025-10.
%  Converted from the old 'Αγαπητέ ασθενή' letter style.
%
%  Copy, rename to "YYYYMMDD - surname - topic.tex", fill in and
%  compile with XeLaTeX. Toggle alternatives by moving the "%".
%  Text in [brackets] is a placeholder: replace or delete it.
% =====================================================================
\documentclass[10pt,greek]{article}
\usepackage[useregional]{datetime2}

%%%%%%%%%%%%%%%%%%%%% Language
\providecommand\DocLang{gr}

%%%%%%%%%%%%%%%%%%%%% Clinic (letterhead logo, footer address)
\providecommand\Clinic{Peristeri}
% \providecommand\Clinic{Mandra}
% \providecommand\Clinic{Bioclinic}
% \providecommand\Clinic{Psychikou}
% \providecommand\Clinic{Metropolitan}
% \providecommand\Clinic{Dunant}
% \providecommand\Clinic{Euroclinic}

%%%%%%%%%%%%%%%%%%%%% Document code for the QR: paste the code issued
%%%%%%%%%%%%%%%%%%%%% by doc.erns.eu (empty = no QR code)
\providecommand\DocCode{}

%%%%%%%%%%%%%%%%%%%%% Patient
\providecommand\PatientPrefix{Ο}
% \providecommand\PatientPrefix{Η κα}
\providecommand\PatientName{Παπαδόπουλος Γιώργος}
\providecommand\PatientAge{45}
% \providecommand\PatientAge{}

%%%%%%%%%%%%%%%%%%%%% Date
\providecommand\DocDate{\DTMdate{2026-01-15}}
% \providecommand\DocDate{\DTMdate{2026-05-04}}

%%%%%%%%%%%%%%%%%%%%% Title
\providecommand\DocTitle{Μετεγχειρητικές Οδηγίες}
% \providecommand\DocTitle{Ενημερωτικό Σημείωμα για τον Ασθενή}

%%%%%%%%%%%%%%%%%%%%% Closing (empty = no signature block)
\providecommand\DocClosing{Ο ιατρός}
% \providecommand\DocClosing{}

%%%%%%%%%%%%%%%%%%%%% Stamp and signature image
\providecommand\DocStamp{no}
% \providecommand\DocStamp{yes}

\input{evrtemplate}

\begin{document}

\begin{PatientInfo}
  \Info{Διάγνωση}{Τιμή 1}
  \Info{Χειρουργείο}{(10/01/2026) Τιμή 2}
  \Info{Ημ. Εξιτηρίου}{10-01-2026}
\end{PatientInfo}

% Short course of the admission (optional, from the old letters):
% Μετά την επέμβαση παραμείνατε για Τιμή 5 ημέρες στη ΜΕΘ, όπου η μετεγχειρητική πορεία ήταν Τιμή 6. Στις Τιμή 7 μεταφερθήκατε στον όροφο, όπου παραμείνατε έως τις Τιμή 8, οπότε και λάβατε εξιτήριο.

Ακολουθούν οδηγίες για την ομαλή σας ανάρρωση:

\section*{Γενικές Οδηγίες}
\begin{enumerate}[label=\arabic*.]
\setlength{\parskip}{6pt}
\item \textbf{Τομή:} Μην βρέξετε το τραύμα (λουτρό/πισίνα) μέχρι την αφαίρεση ραμμάτων ή την επανεξέταση.
\item \textbf{Φροντίδα Τραύματος:} Το τραύμα πρέπει να παραμένει καθαρό και στεγνό. Αλλάξτε επίδεσμο όπως σας έχει υποδειχθεί.
\item \textbf{Αλλαγές:} Αλλαγή τραύματος κάθε 2 ημέρες με αυτοκόλλητα γαζάκια για 7 ημέρες. Μετά μπορεί να χρησιμοποιείται spray τύπου Pulvo47 / Vioplex-T ή άλλο αντίστοιχο, 2 φορές την ημέρα.
\item \textbf{Διατροφή:} Ελαφριά και πλούσια σε πρωτεΐνες και φυτικές ίνες, με επαρκή ενυδάτωση.
\end{enumerate}

\section*{Επανέλεγχοι}
\begin{itemize}
\item Σε 7--10 ημέρες για κλινική επανεξέταση και αφαίρεση ραμμάτων.
\item Σε Τιμή 9 ημέρες λήψη αίματος για αιματολογικές εξετάσεις: γενική αίματος, ουρία, κρεατινίνη, νάτριο, κάλιο, CRP.
% \item Εντός Τιμή 10 εβδομάδων θα ξεκινήσουν συνεδρίες χημειοθεραπείας και ακτινοθεραπείας, σε συνεννόηση με τον ογκολόγο.
\end{itemize}

% Emergency box: not in the old letters. Review before enabling.
% \begin{UrgentBox}
% \item \textbf{Επιληπτική κρίση}.
% \item \textbf{Επιδείνωση του επιπέδου συνείδησης}, έντονη υπνηλία ή σύγχυση.
% \end{UrgentBox}

\textbf{Επικοινωνήστε άμεσα με τον ιατρό εάν εμφανίσετε:}
\begin{itemize}
\item Έντονο ή αυξανόμενο πόνο.
\item Εμέτους, θολή όραση ή διπλωπία.
\item Πυρετό.
\item Ερυθρότητα ή εκροή από την τομή.
\item Αδυναμία ή αιμωδίες στο σώμα που δεν προϋπήρχαν.
\end{itemize}

\newpage
\begin{center}
\LARGE\textbf{Οδηγίες Φαρμακευτικής Αγωγής}
\end{center}
\vspace{0.5cm}

\begin{MedTable}
\DrugEntry{inj Fraxiparine}{2850IU}{1 x 1}{10}{αντιπηκτικές ενέσεις υποδορίως, την ίδια ώρα κάθε ημέρα}
\DrugEntry{tb Klaricid}{500mg}{1 x 2}{7}{αντιβίωση}
\DrugEntry{tb Keppra}{1000mg}{1 x 3}{}{}
\DrugEntry{tb Dexamethasone}{4mg}{1 x 3}{}{μετά το φαγητό}
\DrugEntry{caps Omeprazole}{40mg}{1 x 2}{10}{προ φαγητού}
\DrugEntry{tb Paracetamol}{1000mg}{1 x 3}{10}{}
\DrugEntry{spray Vioplex-T ή Pulvo47}{}{1 x 2}{}{στο τραύμα, και να καλύπτεται μετά}
% ---- Drug library: uncomment the lines you need ---------------------
% -- Analgesics
% \DrugEntry{tb Paracetamol}{1000mg}{1 x 3}{10}{}
% \DrugEntry{tb Apotel}{1000mg}{2 ή 3 την ημέρα}{}{παυσίπονο, μέγιστη δόση 3 δισκία το 24ωρο}
% \DrugEntry{tb Depon}{500mg}{1 ή 2, έως 3 φορές την ημέρα}{}{εφόσον υπάρχει πόνος -- μέγιστο 3g το 24ωρο}
% \DrugEntry{tb Tramal}{50mg}{1 x 2 ή 3}{6}{}
% \DrugEntry{tb Zaldiar}{}{2 ή 3 την ημέρα}{}{εφόσον υπάρχει έντονος πόνος}
% -- NSAIDs
% \DrugEntry{tb Movatec}{7.5mg}{1 x 2}{10}{μετά το φαγητό}
% \DrugEntry{tb Movatec}{15mg}{1 x 1}{1}{μετά το φαγητό}
% \DrugEntry{tb Naprosyn}{500mg}{1 x 2}{6}{μετά το φαγητό}
% -- Steroids (tapers)
% \DrugEntry{tb Medrol}{4mg}{1 x 3\newline 1 x 2\newline 1 x 1\newline STOP}{}{για 5 ημέρες\newline για 5 ημέρες\newline για 5 ημέρες}
% \DrugEntry{tb Medrol}{4mg}{2 x 3\newline 1 x 3\newline 1 x 2\newline 1 x 1\newline STOP}{}{για 3 ημέρες\newline για 3 ημέρες\newline για 3 ημέρες\newline για 3 ημέρες}
% \DrugEntry{syr Dexaton}{2mg/5ml}{10cc x 3}{10}{μετά το φαγητό}
% \DrugEntry{tb Dexamethasone}{4mg}{1 x 3}{}{μετά το φαγητό}
% -- Antibiotics
% \DrugEntry{tb Zinadol}{500mg}{1 x 2}{5}{αντιβίωση}
% \DrugEntry{tb Klaricid}{500mg}{1 x 2}{7}{αντιβίωση}
% \DrugEntry{caps Amoxicillin}{500mg}{1 x 3}{7}{μετά το φαγητό}
% -- Gastroprotection
% \DrugEntry{caps Omeprazole}{40mg}{1 x 2}{10}{προ φαγητού}
% \DrugEntry{caps Omeprazole}{20mg}{1 x 2}{10}{15 λεπτά πριν το φαγητό}
% \DrugEntry{caps Mostrelan}{40mg}{1 x 1}{10}{προ φαγητού}
% -- Anticoagulants
% \DrugEntry{inj Clexane}{4000IU}{1 x 1}{10}{αντιπηκτικές ενέσεις υποδορίως, την ίδια ώρα κάθε ημέρα}
% \DrugEntry{inj Innohep}{0,45ml}{1 x 1}{10}{αντιπηκτικές ενέσεις υποδορίως, την ίδια ώρα κάθε ημέρα}
% \DrugEntry{inj Fraxiparine}{2850IU}{1 x 1}{10}{αντιπηκτικές ενέσεις υποδορίως, την ίδια ώρα κάθε ημέρα}
% -- Neuropathic pain, antiepileptics
% \DrugEntry{tb Lyrica}{50mg}{1 x 1 (βράδυ)}{4}{μετά την 4η ημέρα, 1 x 2}
% \DrugEntry{caps Lyrica}{150mg}{1 x 2}{}{πρωί -- βράδυ, με τροφή. Μην διακόπτεται απότομα}
% \DrugEntry{tb Neurontin}{400mg}{2 x 2}{}{}
% \DrugEntry{tb Tegretol}{200mg}{1 x 3}{}{με τροφή, σταθερές ώρες. \textbf{Μην διακόπτεται απότομα}}
% \DrugEntry{caps Cymbalta}{60mg}{1 x 2}{}{}
% \DrugEntry{tb Keppra}{1000mg}{1 x 3}{}{}
% \DrugEntry{tb Depakine}{400mg}{1 x 3}{}{θα γίνει προοδευτική μείωση σε συνεννόηση με τον θεράποντα ιατρό}
% -- Other
% \DrugEntry{sir Duphalac}{5ml}{1 ή 2 την ημέρα}{}{καθαρτικό, ανάλογα με τη λειτουργία του εντέρου -- στόχος: μαλακές κενώσεις χωρίς προσπάθεια}
% \DrugEntry{tb Xanax}{0,25mg}{1 x 2-3}{}{μέχρι 3 την ημέρα επί εκνευρισμού}
% \DrugEntry{tb Escitalopram}{10mg}{1 x 1}{}{}
% \DrugEntry{spray Vioplex-T ή Pulvo47}{}{1 x 2}{}{στο τραύμα, και να καλύπτεται μετά}
% ---------------------------------------------------------------------
\end{MedTable}

% Notes under the table (uncomment the ones you need)
\textbf{Σημείωση:} Μην διακόψετε και μην μειώσετε μόνος/μόνη σας τα Τιμή 11. Οι δόσεις θα μειωθούν σταδιακά κατόπιν ιατρικής οδηγίας.
% \textbf{Σημείωση:} Μην λαμβάνετε αντιφλεγμονώδη φάρμακα (π.χ.\ Voltaren, Brufen, Ponstan, Naprosyn) χωρίς ιατρική οδηγία.

\end{document}
